
\begin{table}[htbp]
   \caption{\label{tab:robust_state} Robustness: State-Level Mobilization Effect on Female LFP Change}
   \centering
   \begin{tabular}{lcccc}
      \tabularnewline \midrule \midrule
      Dependent Variable: & \multicolumn{4}{c}{$\Delta$ Female LFP}\\
                               & Baseline       & Unweighted & Non-South & Trimmed \\   
      Model:                   & (1)            & (2)        & (3)       & (4)\\  
      \midrule
      \emph{Variables}\\
      Mobilization Rate (std.) & -0.0073$^{**}$ & -0.0015    & -0.0056   & -0.0013\\   
                               & (0.0027)       & (0.0036)   & (0.0040)  & (0.0033)\\   
      \midrule
      \emph{Fit statistics}\\
      Observations             & 49             & 49         & 33        & 39\\  
      R$^2$                    & 0.67586        & 0.38254    & 0.60773   & 0.63769\\  
      Adjusted R$^2$           & 0.62052        & 0.29433    & 0.51720   & 0.56975\\  
      \midrule \midrule
      \multicolumn{5}{l}{\emph{IID standard-errors in parentheses}}\\
      \multicolumn{5}{l}{\emph{Signif. Codes: ***: 0.01, **: 0.05, *: 0.1}}\\
   \end{tabular}
   
   \par \raggedright 
   Dependent variable: change in female LFP 1940--1950. Column (1): baseline with 1940 controls, population-weighted. Column (2): unweighted. Column (3): excludes 16 former Confederate states plus border states. Column (4): trims top and bottom 10\% of mobilization distribution. All specifications include 1940 controls (percent urban, Black, farm, mean education, mean age).
\end{table}


